\documentclass[11pt,a4paper]{article}
\usepackage[T1]{fontenc}
\usepackage[utf8]{inputenc}
\usepackage[english]{babel}
\usepackage{amsmath,amssymb}
\usepackage[margin=2cm]{geometry}
\setlength{\parindent}{0pt}
\setlength{\emergencystretch}{3em}
\allowdisplaybreaks

\title{About \texttt{book\_partial\_derivatives.tex}}
\author{}
\date{}
\begin{document}
\maketitle
This note explains what the test data in \texttt{book\_partial\_derivatives.tex} is used for. The file is read by \texttt{test\_multivariable\_calculus.c}.

\begin{enumerate}
\item \textbf{Purpose.} The file holds 71 results printed in collections of solved exercises of a second course of calculus: partial derivatives, gradients and Hessian matrices, as formulas or at a point. Many points are the critical points of the exercises on maxima and minima, with the Hessian matrix that the exercise then uses to say what kind of point it is; saying it is not part of this file.

\item \textbf{What it exercises.} The same as \texttt{partial\_derivatives.tex}, on the functions that the exercises actually use: quotients, logarithms and roots with a domain that is a region of the plane (items 25 and 26 are defined in no sample point of chapter 27), a power $x^{y}$, powers with a rational exponent, three variables with a $3 \times 3$ Hessian matrix, points with fractions, with $\sqrt{3}$, with $e^{-1/4}$ and with $\frac{1}{e}$ among the coordinates or the values.

\item \textbf{Expected result.} Every statement is proved with a written proof; 69 carry an accepted certificate. Item 11 (the second derivatives of $x^{y}$) and item 57 (a gradient at a point with $\sqrt{3}$) are proved without it. The other checks are those of \texttt{partial\_derivatives.tex}.

\item \textbf{Sources.} The PDF files are in \texttt{photos/funzioni\_due\_variabili/}; the pages and the check of each result by a third computation are in \texttt{photos/dati\_esercizi/analisi\_funzioni\_due\_variabili.md}, section 5. Items 1--2: CTU Prague, Habala, n.~2 (\texttt{ctu\_habala\_ema2r2.pdf}, p.~1, 3). Items 3--6, 37--39, 57--58: Pavia, Vitali (\texttt{unipv\_vitali\_am2\_derivazione1.pdf}, p.~1, 3--5). Items 7--13, 40--41, 52--54: Verona, Eleuteri (\texttt{univr\_eleuteri\_eserciziario\_analisi2.pdf}, 3.4.2--3.4.7, 3.6.6, 3.6.7, 3.8.1, 5.1.2, 5.1.9, 5.1.11). Items 14--17, 55--56, 60: Tor Vergata, Damascelli (\texttt{uniroma2\_damascelli\_esempi\_analisi2\_chimica\_wayback.pdf}, p.~1--2, 7--9, 11). Items 18--21, 42--44: UBC, Laba and Marshall, 2.4.1, 2.4.2, 2.5.1, 2.5.2, 3.3.3. Items 22--23: Sapienza, Bersani, n.~8. Items 24--28, 32--33: Trieste (\texttt{units\_piuvar\_svolti.pdf}, n.~4, 6, 12). Item 29: Trieste, Obersnel, lesson 5. Items 30--31, 59: Verona, Pagliarini, n.~6, 11. Items 34--36: Habala n.~3, 4, Vitali n.~3c. Item 45: Insubria, n.~5. Item 46: Eleuteri, notes on differentiability, 1.8.6. Item 47: Tor Vergata, Perfetti, 6.7. Items 48--49: Bologna (\texttt{unibo\_es\_funz\_2var\_1.pdf}). Items 50--51: MIT 18.02SC, second derivative test. Items 61--71 are from \texttt{photos/estremi\_vincolati/} (\texttt{photos/dati\_esercizi/analisi\_estremi\_vincolati.md}, section 4c): items 61--68 from Modena, Magni and Vezzali (\texttt{unimore\_magni\_fascicolo4\_esercizi.pdf}, 14.12, 14.22, 14.23, p.~5--13), where the determinant of the Hessian matrix at a point is written as an expression of second derivatives; items 69--71 from an MIT exam and two Berkeley Math 53 worksheets.

The printed results that are wrong are in \texttt{false\_book\_partial\_derivatives.tex}, and here there is the right one: items 21 (UBC 2.5.2b), 30 (Pagliarini 11), 40 (Eleuteri 3.4.2), 52 (Eleuteri 5.1.2), 54 (Eleuteri 5.1.11). All 71 were confirmed again by the independent oracle of \texttt{tools/fuzz/fuzz\_multivariable.py} (mode \texttt{check}).

\item \textbf{How to run it.} From the top folder of the repository:
\begin{verbatim}
cc -std=c11 -O2 -Wall -Wextra -o /tmp/test_multivariable_calculus \
    tests/multivariable_calculus/test_multivariable_calculus.c && \
    /tmp/test_multivariable_calculus
\end{verbatim}
\end{enumerate}
\end{document}
